% High-level overview of the IVF search with per-list query blocks and two passes
% (ann_ivf_ttaiss), drawn as a host/device timeline.
% Requires: \usepackage{tikz,amsmath,xcolor}
%           \usetikzlibrary{calc,arrows.meta}
\providecolor{hostc}{RGB}{214,228,240}
\providecolor{devc}{RGB}{226,236,214}

% Background (optional; can be moved to the section text).
\paragraph{Background: IVF search.} The goal of approximate nearest-neighbour
(ANN) search is to find, for every query vector, the \(k\) dataset vectors most
similar to it; we use \(k = 10\) and cosine similarity (inner product of
L2-normalised vectors). An inverted-file (IVF) index avoids comparing each query
with the whole dataset. K-means splits the dataset into \(N_{\mathrm{list}}\)
clusters; each cluster has a centroid and a list of the vectors closest to it.
A query is first compared with the \(N_{\mathrm{list}}\) centroids (coarse
search). Then only the vectors of its \(N_{\mathrm{probe}}\) closest clusters are
scanned (fine search), and the best \(k\) of them are returned. A larger
\(N_{\mathrm{probe}}\) scans more vectors and gives a higher recall (the fraction
of the true 10 nearest neighbours found). Our experiments use GloVe-100
(1{,}183{,}514 vectors of dimension 100) with \(N_{\mathrm{list}} = 512\) and
\(N_{\mathrm{probe}} = 64\), and search batches of 10{,}000 queries. Each query
then scans 64 of the 512 clusters, one eighth of them.

\begin{figure}[htbp]
\centering
\resizebox{0.95\textwidth}{!}{%
\begin{tikzpicture}[font=\small,
  hostbox/.style={draw,rounded corners=4pt,fill=hostc,align=left,text width=5.0cm,inner sep=5pt},
  devbox/.style={draw,rounded corners=4pt,fill=devc,thick,align=left,text width=5.6cm,inner sep=5pt},
  arr/.style={-{Stealth[length=2mm]},thick},
  lbl/.style={font=\scriptsize,fill=white,inner sep=1pt,align=center}]
% lane centres
\coordinate (hx) at (0,0);
\coordinate (dx) at (8.2,0);
% context
\node[draw,rounded corners=4pt,fill=gray!8,align=left,text width=13.6cm,inner sep=6pt,anchor=south]
  (ctx) at (4.1,1.3) {%
  \textbf{Goal:} for every query, the \(k=10\) most similar dataset vectors (cosine similarity).\\
  \textbf{IVF index:} k-means splits the dataset into \(N_{\mathrm{list}}\) clusters (centroid + list of vectors).\\
  \textbf{Search:} compare the query with the centroids, scan only the vectors of its \(N_{\mathrm{probe}}\)
  closest clusters, keep the best \(k\).\\
  \textbf{Here:} GloVe-100 (1.18\,M vectors, \(d=100\)), \(N_{\mathrm{list}}=512\), \(N_{\mathrm{probe}}=64\),
  batches of 10{,}000 queries.};
\node[font=\bfseries] at (0,0.8) {Host (CPU)};
\node[font=\bfseries] at (8.2,0.8) {Device (Tenstorrent, 64 cores)};

% 1. index, once
\node[hostbox,anchor=north] (h1) at (hx) {%
  \textbf{1. Once:} upload the index\\
  (centroids, cluster vectors + IDs)};
\node[devbox,anchor=north] (d1) at (dx |- h1.north) {%
  Index kept in device memory\\for all searches};
\draw[arr] (h1.east) -- (d1.west);

% 2. queries
\node[hostbox,anchor=north] (h2) at ($(hx |- h1.south)+(0,-0.6)$) {%
  \textbf{2. Per search:} upload the queries};

% 3. coarse search
\node[devbox,anchor=north] (d3) at ($(dx |- h2.south)+(0,-0.3)$) {%
  \textbf{3. Coarse search}\\
  queries \(\times\) centroids; keep the closest\\
  centroids of each query (top-64)};
\draw[arr] (h2.east) -- node[lbl,above,sloped]{queries} (d3.west);

% 4. plan pass 1
\node[hostbox,anchor=north] (h4) at ($(hx |- d3.south)+(0,-0.3)$) {%
  \textbf{4. Plan}\\[2pt]
  \textbullet\ keep the \(N_{\mathrm{probe}}\) closest clusters of each query\\
  \textbullet\ pass 1 = its \(N_{\mathrm{probe}}/8\) closest clusters (8 for\\
  \hspace*{0.8em}\(N_{\mathrm{probe}}=64\)), pass 2 = the others\\
  \textbullet\ for every cluster, put the queries that\\
  \hspace*{0.8em}probe it in blocks of \(\le 32\)\\[2pt]
  \emph{Example:} q1 and q2 both probe \(C_3\)\\
  \(\rightarrow\) \(C_3\) gets the block \(\{q_1, q_2, \ldots\}\)};
\draw[arr] (d3.west) -- node[lbl,above,sloped]{closest clusters} (h4.east);

% 5. pass 1
\node[devbox,anchor=north] (d5) at ($(dx |- h4.south)+(0,-0.3)$) {%
  \textbf{5. Pass 1} (each query's 8 closest clusters)\\[2pt]
  for every cluster: read its vectors once for\\
  up to 32 of its query blocks, multiply them\\
  with each block \(\rightarrow\) top-32 of each query\\
  \emph{inside} that cluster\\[2pt]
  merge the 8 lists of each query \(\rightarrow\) saved:\\
  \emph{pass-1 result} = top-32 over its 8 clusters\\
  bound \(T_q\) = 10th best score of that result};
\draw[arr] (h4.east) -- node[lbl,above,sloped]{blocks (query IDs)} (d5.west);

% 6. plan pass 2, in parallel with pass 1
\node[hostbox,anchor=north] (h6) at (hx |- d5.north) {%
  \textbf{6. Plan pass 2} for the other clusters\\
  (while the device runs pass 1)};

% 7. pass 2
\node[devbox,anchor=north] (d7) at ($(dx |- d5.south)+(0,-0.6)$) {%
  \textbf{7. Pass 2} (the other clusters)\\[2pt]
  same as pass 1, but for each 32 vectors \(\times\)\\
  query block the scores are computed and\\
  the top-32 update is skipped when no score\\
  beats the bound \(T_q\) of its query\\[2pt]
  saved: one top-32 list per (query, cluster),\\
  \emph{not merged yet}; empty if all skipped};
\draw[arr] (d5.south) -- node[lbl,right]{\(T_q\)} (d7.north);
\draw[arr] (h6.east) -- node[lbl,above,sloped]{pass-2 blocks} (d7.west);

% 8. final top-k
\node[devbox,anchor=north] (d8) at ($(dx |- d7.south)+(0,-0.6)$) {%
  \textbf{8. Final top-32 per query}\\[2pt]
  \texttt{for each query q (32 at a time):}\\
  \texttt{\ \ best = pass-1 result of q}\\
  \texttt{\ \ for each pass-2 cluster j of q:}\\
  \texttt{\ \ \ \ best = top32(best} \(\cup\) \texttt{list of q in j)}\\[2pt]
  \(\rightarrow\) top-32 of q over all its \(N_{\mathrm{probe}}\) clusters};
\draw[arr] (d7.south) -- (d8.north);

% 9. results
\node[hostbox,anchor=north] (h9) at ($(hx |- d8.south)+(0,-0.3)$) {%
  \textbf{9. Results}\\
  read the top-32 of each query;\\
  optional FP32 re-ranking;\\
  return the top-10};
\draw[arr] (d8.west) -- node[lbl,above,sloped]{top-32} (h9.east);

% lane separator
\draw[dashed,gray] (4.1,1.0) -- (4.1,0 |- h9.south);
\end{tikzpicture}}
\caption{Overview of one search for the \(k=10\) nearest neighbours of every
query in a batch. The index is uploaded once. For each batch
of queries, the device finds the closest clusters of every query (coarse
search); the host groups the queries by cluster, so that each cluster is read
once and multiplied with blocks of the queries that need it. Pass 1 searches
each query's closest clusters and gives a bound \(T_q\); pass 2 searches the
remaining clusters and skips the blocks that cannot beat that bound. A final
top-32 merge combines all results per query; the host returns the top-10
(after an optional FP32 re-ranking). With re-ranking, \(T_q\) is the 32nd best
score instead of the 10th. Below \(N_{\mathrm{probe}}=24\) there is a single pass
over all \(N_{\mathrm{probe}}\) clusters, followed by the final merge.}
\label{fig:gf-overview}
\end{figure}
